\chapter{Mean-Dominated Quantization Noise and Softmax Amplification}
\label{ch:softmax_amplification_full}

This chapter explains a concrete and practically important phenomenon in low-bit
Transformer training: \emph{even when stochastic rounding is unbiased, quantization
noise can be amplified by softmax and lead to instability}. The central message is
that the amplification is not mysterious; it follows from a precise causal chain.
A coherent mean bias in the upstream activations creates extreme elementwise values,
which in turn \emph{set the quantization scale}. This ``scale hijacking'' compresses
the long-tail semantic variations into a tiny portion of the quantization grid.
Stochastic rounding remains unbiased, but the resulting variance becomes large
\emph{relative to the long-tail signal}. Softmax then propagates this variance through
a dense Jacobian coupling term and, additionally, converts variance into systematic
drift through curvature (Hessian) terms.

We develop this argument in a slow, step-by-step manner. At each step we keep track
of what is (i) an algebraic identity, (ii) an approximation (Taylor expansion), and
(iii) a modeling assumption (noise model, tail behavior). The point is not to
``prove softmax is bad,'' but to understand the mechanism well enough to design
stable low-bit systems.

\section{Preliminaries: what softmax acts on in Transformers}

In attention, softmax is applied to a vector of scores over \emph{positions}.
Fix a query token index $i$ and let the key sequence length be $n$ (typically $n=L$).
The pre-softmax scores for query $i$ are
\begin{equation}
z \in \mathbb{R}^{n},\qquad
z_j = \frac{\langle q_i, k_j\rangle}{\sqrt{d}},
\end{equation}
and the attention weights are
\begin{equation}
y = \mathrm{softmax}(z/\tau)\in\mathbb{R}^n,\qquad
y_j = \frac{\exp(z_j/\tau)}{\sum_{\ell=1}^{n}\exp(z_\ell/\tau)},
\end{equation}
where $\tau>0$ is an effective temperature (explicit or implicit).

\paragraph{Why we study the score vector $z$.}
Low-bit failures in attention often arise from quantizing the scores (or intermediate
quantities that determine them). The important point is that the softmax output
depends sensitively on the \emph{relative} structure of $z$, especially when $y$
becomes sharp (peaky). Our goal is to understand how quantization noise in $z$
propagates through softmax.

\paragraph{Checkpoint (critical thinking).}
Before doing any math, ask: \emph{what controls training stability?} In low-bit
training, it is rarely the average energy; it is often the \emph{extremes} that set
scales and dominate variance. This chapter formalizes that intuition for the softmax
block.

\section*{Why a Small Mean Becomes Dominant in High Dimension}

Before analyzing quantization and softmax, we must first explain a simple but
crucial geometric fact: in high-dimensional spaces, a coherent mean bias can be
small in each coordinate yet large in total magnitude.

Let $\mu \in \mathbb{R}^H$ denote the mean vector of an activation (or logit)
distribution, where $H$ is the hidden dimension (typically $H \ge 1024$).
Suppose that the mean bias is \emph{coherent}, meaning that it shifts many
coordinates in a consistent direction.
If the typical coordinate-wise mean magnitude is $\bar{\mu}$, i.e.,
$|\mu_j| \approx \bar{\mu}$ for many $j$, then the Euclidean norm satisfies

\begin{equation}
\|\mu\|_2
=
\sqrt{\sum_{j=1}^{H} \mu_j^2}
\approx
\sqrt{H}\,\bar{\mu}.
\label{eq:mu_scaling}
\end{equation}

This $\sqrt{H}$ factor is the high-dimensional magnifier.
For example, when $H=4096$, we have $\sqrt{H}=64$.
Thus a coordinate-wise mean shift as small as $\bar{\mu}=0.01$ produces a global
mean magnitude $\|\mu\|_2 \approx 0.64$, which is no longer small relative to the
scale of typical fluctuations.

The contrast with unstructured noise is equally important.
If $x\in\mathbb{R}^H$ is mean-zero with per-coordinate variance $\sigma^2$,
then the typical coordinate magnitude is $O(\sigma)$, and the maximum coordinate
magnitude grows only as $O(\sigma\sqrt{2\log H})$ by standard concentration.
In other words, coherent bias grows with $\sqrt{H}$, while the largest random
fluctuation grows only with $\sqrt{\log H}$.

This difference in scaling explains why coherent mean components dominate extreme
elementwise behavior.
Even if the mean accounts for only a modest fraction of total variance,
its coherence across coordinates makes it disproportionately likely to appear in
the largest entries.
Since low-bit quantization and adaptive optimizers are controlled by such
extremes (blockwise maxima and squared magnitudes), mean bias becomes the
first-order driver of numerical instability.

With this geometric principle in mind, we can now analyze the full causal chain:
mean-dominated outliers determine the quantization scale, quantization injects
variance into the long tail, and softmax propagates and amplifies that variance.

\paragraph{A concrete numerical example.}
To make \eqref{eq:mu_scaling} tangible, consider a typical LLM hidden size
$H=4096$ (so $\sqrt{H}=64$). Suppose the mean bias is small per coordinate,
with typical magnitude $\bar{\mu}=0.01$. Then
\[
\|\mu\|_2 \approx \sqrt{H}\,\bar{\mu} = 64 \times 0.01 = 0.64.
\]
So even though the bias is only $0.01$ in each coordinate, its coherent
accumulation across thousands of dimensions produces a global shift of order
$0.64$, which is not negligible.

Now compare this to random mean-zero fluctuations.
Assume a mean-zero noise vector with per-coordinate standard deviation
$\sigma=0.1$. A typical coordinate has magnitude on the order of $0.1$, but the
largest coordinate among $H$ dimensions grows only logarithmically:
\[
\max_{1\le j\le H} |x_j|
\;\approx\;
\sigma \sqrt{2\log H}
\;=\;
0.1 \sqrt{2\log 4096}.
\]
Since $\log 4096 \approx 8.32$ (natural log), this gives
\[
\max_j |x_j|
\approx
0.1 \sqrt{2\times 8.32}
=
0.1 \sqrt{16.64}
\approx
0.1 \times 4.08
\approx
0.41.
\]
Thus, in this example, the coherent mean produces a global magnitude
$\|\mu\|_2 \approx 0.64$, while even the \emph{largest} random coordinate is
only about $0.41$.
This illustrates the key high-dimensional contrast:
coherent bias scales as $\sqrt{H}$, whereas random extremes scale only as
$\sqrt{\log H}$.

For low-bit quantization, this difference is decisive.
Blockwise scaling is set by extreme magnitudes, and a coherent mean shifts
many coordinates in the same direction.
Even if each coordinate shift is small, the induced global drift is large enough
to dominate the extreme statistics that control quantization scale and
second-moment accumulation.

\section{Stage I: scale hijacking --- why the quantizer grid is set by bias}

We model the score vector as
\begin{equation}
z = m + r,
\label{eq:z_m_r_ch1}
\end{equation}
where $m$ denotes a coherent component (bias-driven contribution) and $r$ denotes
the residual semantic structure (spike+tail). The crucial empirical fact established
earlier is that coherent bias often dominates \emph{elementwise} extremes of upstream
activations. When these activations are linearly transformed into scores, the same
pattern can appear in $z$: the extreme values of $z$ can be dominated by $m$.

\subsection{Block quantization and step size}

Consider symmetric $b$-bit uniform quantization applied blockwise. Let $s$ be the
scale chosen for the block and define the step size
\begin{equation}
\Delta = \frac{s}{2^{b-1}-1}.
\label{eq:delta_def_ch1}
\end{equation}
In many implementations, $s$ is chosen as either the maximum magnitude or a high
percentile of magnitudes within the block. For pedagogy, assume the simplest and
most extreme choice:
\begin{equation}
s = \|z\|_\infty = \max_{1\le j\le n}|z_j|.
\label{eq:scale_inf}
\end{equation}

\subsection{Why extremes decide the grid}

Substituting $z=m+r$, the scale satisfies
\begin{equation}
\|z\|_\infty = \|m+r\|_\infty \ge \|m\|_\infty - \|r\|_\infty.
\end{equation}
If $\|m\|_\infty \gg \|r\|_\infty$, then $\|z\|_\infty$ is essentially set by $\|m\|_\infty$.
This is already enough to explain the phenomenon at an engineering level.

To connect this to high-dimensional probability, assume the residual entries $r_j$
are sub-Gaussian with parameter $\sigma^2$, meaning
\begin{equation}
\Pr(|r_j|>t)\le 2\exp\!\left(-\frac{c t^2}{\sigma^2}\right)
\quad\text{for some }c>0.
\label{eq:subgauss_tail}
\end{equation}
Then standard extreme-value scaling gives
\begin{equation}
\|r\|_\infty \approx \sigma\sqrt{2\log n}
\quad\text{with high probability}.
\label{eq:residual_max_scale}
\end{equation}
Therefore, whenever the coherent component satisfies
\begin{equation}
\|m\|_\infty \gg \sigma\sqrt{2\log n},
\label{eq:hijack_condition}
\end{equation}
the block scale $s=\|z\|_\infty$ is mean-dominated with high probability.
We call this phenomenon \emph{scale hijacking}.

\paragraph{Interpretation.}
Scale hijacking means the quantizer spends its limited codebook to represent the
coherent component $m$ accurately (because $m$ sets the dynamic range), while the
residual semantic structure $r$ is forced into the small remaining numerical range.

\paragraph{Checkpoint (critical thinking).}
Students often assume that ``unbiased noise is harmless.'' This chapter shows that
the \emph{variance} of unbiased noise is what matters, and the variance is controlled
by the quantizer step size $\Delta$, which is controlled by \emph{outliers}. If
outliers are mean-dominated, the quantization grid becomes biased toward representing
the mean rather than the semantics.

\section{Stage II: variance inflation --- why unbiased rounding still causes trouble}

\subsection{Stochastic rounding as additive noise}

Model stochastic rounding as
\begin{equation}
\hat z = z + \varepsilon,
\qquad
\mathbb{E}[\varepsilon_j\mid z_j]=0.
\label{eq:sr_noise_model}
\end{equation}
For uniform quantization, the error is bounded by half a step, hence
\begin{equation}
\mathrm{Var}(\varepsilon_j\mid z_j)\le \frac{\Delta^2}{4}.
\label{eq:sr_var_bound}
\end{equation}
A common approximation is
\begin{equation}
\mathrm{Var}(\varepsilon_j\mid z_j)\approx \frac{\Delta^2}{12},
\label{eq:sr_var_approx}
\end{equation}
which holds if the rounding error is close to uniform within a bin.

\subsection{Tail SNR collapse}

Consider a coordinate $j$ where the semantic residual magnitude satisfies $|r_j|\ll \Delta$.
Then the quantized residual $\hat r_j=r_j+\varepsilon_j$ is dominated by $\varepsilon_j$,
and the effective signal-to-noise ratio obeys
\begin{equation}
\mathrm{SNR}_j
\triangleq
\frac{|r_j|}{\sqrt{\mathrm{Var}(\varepsilon_j)}}
\lesssim
\frac{|r_j|}{\Delta}.
\label{eq:snr_tail}
\end{equation}
Under scale hijacking, $\Delta$ is proportional to $s\approx \|m\|_\infty$, so
\begin{equation}
\mathrm{SNR}_j \lesssim \frac{|r_j|}{\|m\|_\infty}.
\label{eq:snr_mean_control}
\end{equation}
This is the key point: stochastic rounding is unbiased, but the variance can be
enormous relative to the long-tail semantic signal.

\paragraph{Interpretation.}
Mean bias does not need to ``directly break'' learning. It suffices that it forces
the quantizer step size to be large. Then the quantized tail components become noisy,
and that noise is fed into softmax.

\section{Stage III: dense coupling amplification --- how softmax spreads and amplifies noise}

Softmax is nonlinear and dense: every output coordinate depends on every input coordinate.
This is what turns localized variance into global instability.

\subsection{Taylor expansion framework}

Let $y=\mathrm{softmax}(z/\tau)$ and $\hat y=\mathrm{softmax}(\hat z/\tau)$.
A second-order Taylor expansion yields
\begin{equation}
\hat y_i
=
y_i
+
\sum_{j=1}^{n} J_{ij}(z)\,\varepsilon_j
+
\frac12 \sum_{j=1}^{n}\sum_{k=1}^{n} H_{i,jk}(z)\,\varepsilon_j\varepsilon_k
+
O(\|\varepsilon\|^3),
\label{eq:taylor_softmax_ch1}
\end{equation}
where $J$ is the Jacobian and $H$ is the Hessian tensor of softmax.

\subsection{Jacobian: first-order variance propagation}

The softmax Jacobian is
\begin{equation}
J(z)
=
\frac{1}{\tau}\left(\mathrm{diag}(y)-yy^\top\right),
\qquad
J_{ij}(z)=\frac{1}{\tau}y_i(\delta_{ij}-y_j).
\label{eq:softmax_jac_ch1}
\end{equation}
To first order, $\Delta y \approx J\varepsilon$, hence
\begin{equation}
\mathrm{Cov}(\Delta y)\approx J\,\Sigma_\varepsilon\,J^\top,
\label{eq:cov_prop_ch1}
\end{equation}
where $\Sigma_\varepsilon=\mathrm{Cov}(\varepsilon)$.

A crucial implication is that the covariance is not diagonal even if $\Sigma_\varepsilon$
is diagonal, because $J$ is dense. This is the \emph{dense coupling} effect.

\subsection{Head-to-tail broadcasting in the peaky regime}

When attention is sharp, one coordinate $h$ dominates: $y_h\approx 1$ and $y_i\ll 1$ for $i\neq h$.
In this regime, for $i\neq h$ we have
\begin{equation}
J_{ih}(z)= -\frac{1}{\tau}y_i y_h \approx -\frac{1}{\tau}y_i.
\end{equation}
If $\Sigma_\varepsilon$ is diagonal with entrywise variances $\sigma_j^2$, then
\begin{equation}
\mathrm{Var}(\Delta y_i)
=
\sum_{j=1}^{n}J_{ij}^2\sigma_j^2
\approx
J_{ih}^2\sigma_h^2
\approx
\frac{1}{\tau^2}y_i^2\sigma_h^2.
\label{eq:head_tail_broadcast_ch1}
\end{equation}
Thus the variance of the \emph{dominant head coordinate} is broadcast into all tail coordinates.
If scale hijacking makes $\sigma_h^2$ large, then all tail probabilities inherit that noise.

\paragraph{Interpretation.}
This is precisely why softmax is dangerous under mean-dominated quantization noise:
even if only a few coordinates have large quantization variance, softmax spreads it globally.

\subsection{Hessian: second-order systematic drift}

Taking expectation of \eqref{eq:taylor_softmax_ch1} under unbiased noise
$\mathbb{E}[\varepsilon]=0$ removes the first-order term but leaves the second-order term:
\begin{equation}
\mathbb{E}[\hat y_i-y_i]
\approx
\frac12 \sum_{j,k} H_{i,jk}(z)\,\Sigma_{\varepsilon,jk}.
\label{eq:hessian_drift_ch1}
\end{equation}
Since $\Sigma_\varepsilon$ scales as $\Delta^2$ and $\Delta$ is mean-controlled under scale hijacking,
the systematic drift scales quadratically with the mean-dominated scale.
Thus unbiased rounding can still induce a predictable shift in attention weights.

\section{Conclusion of the Amplification Mechanism}

Putting the pieces together, mean bias produces instability through a fully specified chain.
Mean-dominated outliers determine the quantization step size $\Delta$ (scale hijacking).
Large $\Delta$ induces large quantization variance, which disproportionately corrupts long-tail semantic coordinates (variance inflation).
Softmax then propagates and amplifies this variance through a dense Jacobian coupling term (head-to-tail broadcasting), and additionally converts variance into systematic drift through curvature (Hessian) terms.
This explains why mean removal is a first-order stabilization mechanism in low-bit attention: it prevents the mean component from setting $\Delta$ in the first place.


\chapter{Mitigating Softmax Error Amplification via Mean-Bias Removal}
\label{ch:softmax_mitigate_mean}

We now derive how mean-bias removal mitigates quantization-induced error
propagation and amplification through softmax.
The key principle is simple but often misunderstood:
softmax itself is translation-invariant, but the quantizer is not.
Thus mean subtraction does not change the ideal (full-precision) softmax output,
but it can shrink the quantization step size and therefore suppress both
Jacobian variance amplification and Hessian-induced drift.

\section{Softmax Invariance and Quantization Non-Invariance}

Softmax satisfies
\begin{equation}
\mathrm{softmax}\big((z+c\mathbf{1})/\tau\big)=\mathrm{softmax}(z/\tau)
\qquad \forall c\in\mathbb{R}.
\end{equation}
Therefore, centering logits by subtracting their average
\begin{equation}
z_c \triangleq z - \bar z\,\mathbf{1},
\qquad
\bar z = \frac{1}{n}\mathbf{1}^\top z,
\end{equation}
preserves the exact attention distribution in full precision:
\begin{equation}
\mathrm{softmax}(z_c/\tau)=\mathrm{softmax}(z/\tau).
\end{equation}

Quantization, however, depends on magnitude through the scale $s$ and step size $\Delta$.
If $z=m+r$ and $s(z)\approx \|m\|_\infty$, then $\Delta$ is mean-controlled.
After centering, the coherent component is removed, and the scale becomes
residual-controlled: $s(z_c)\approx \|r\|_\infty$.

\section{Quadratic Suppression of Noise Covariance}

Define the scale reduction factor
\begin{equation}
\alpha \triangleq \frac{s(z_c)}{s(z)} \approx \frac{\|r\|_\infty}{\|m\|_\infty}.
\end{equation}
Because $\Delta$ scales linearly with $s$,
quantization variance scales quadratically:
\begin{equation}
\Sigma_\varepsilon(z_c) \approx \alpha^2 \Sigma_\varepsilon(z).
\end{equation}

Thus mean-bias removal suppresses injected quantization noise by $\alpha^2$.

\section{Quadratic Suppression of Jacobian Variance Amplification}

Since $y$ is unchanged in full precision by centering, the Jacobian $J$ is unchanged as well.
Therefore,
\begin{equation}
\mathrm{Cov}(\Delta y\mid z_c)
\approx
J\,\Sigma_\varepsilon(z_c)\,J^\top
\approx
\alpha^2\,J\,\Sigma_\varepsilon(z)\,J^\top.
\end{equation}
Hence Jacobian-propagated variance is reduced by a factor $\alpha^2$.

\section{Quadratic Suppression of Hessian Drift}

Similarly, because the Hessian depends only on $y$, it is unchanged under exact centering.
Thus,
\begin{equation}
\mathbb{E}[\hat y-y\mid z_c]
\approx
\frac12 \sum_{j,k} H_{i,jk}(z)\,\Sigma_{\varepsilon,jk}(z_c)
\approx
\alpha^2\,\mathbb{E}[\hat y-y\mid z].
\end{equation}
Therefore mean-bias removal suppresses curvature-induced drift by the same quadratic factor.

\section{Takeaway}

Mean subtraction does not alter the ideal attention distribution,
but it prevents mean-dominated outliers from hijacking the quantization scale.
By shrinking $\Delta$, it reduces injected variance quadratically,
thereby suppressing both Jacobian variance broadcasting and Hessian drift.
This is the fundamental reason mean removal is a first-order stabilizer for low-bit attention.